\documentclass{article}
\usepackage[T1]{fontenc}
\usepackage{lmodern}
\usepackage{graphicx}
\usepackage{amsmath,amssymb}
\usepackage[a4paper,margin=2cm]{geometry}
\usepackage[absolute,overlay]{textpos}
\setlength{\TPHorizModule}{1mm}
\setlength{\TPVertModule}{1mm}
\usepackage[indonesian]{babel}
\usepackage{hyperref}
\setlength{\emergencystretch}{2em}
% unit: Halaman 23

\begin{document}

% === Halaman 23 (python/native) ===

Kriptanalisis Knapsack

• Sayangnya, algoritma knapsack dinyatakan sudah tidak aman, karena 

knapsack dapat dipecahkan oleh pasangan  kriptografer  Adi Shamir. 

• Shamir merumuskan transformasi yang memungkinkan mereka 

merekonstruksi superincreasing knapsack dari normal knapsack. dalam 

waktu polynomial.

• Baca makalahnya: A polynomial time algorithm for breaking the basic 

Merkle-Hellman cryptosystem 

   (Published in: 23rd Annual Symposium on Foundations of Computer 

Science  (sfcs 1982) )

23

\end{document}
